% исправления 27.05.2005
\chapter{Языки первого порядка}
        \label{first-order-languages}%
Помимо логических связок, в математических рассуждениях часто
встречаются \emph{кванторы}%
\index{Квантор|defin}
\лк для
любого\пк\ ($\forall$) и \лк существует\пк\ ($\exists$).
Например, определение непрерывности начинается словами \лк для
любого положительного $\varepsilon$ найдётся положительное
$\delta$, для которого\,$\dots$\пк. А одна из аксиом теории групп
(существование обратного элемента) записывается так:
$\forall x \exists y ((xy\hm=1)\land(yx\hm=1))$.

Можно сформулировать различные логические законы, включающие в
себя кванторы. Например, высказывание \лк существует такое $x$,
что $A$\пк\ (где $A$\т некоторое свойство объекта $x$) логически
эквивалентно высказыванию \лк не для всех $x$ верно $\lnot
A$\пк.

Мы будем записывать такого рода законы с помощью формул, дадим
определение истинности формул (при данной интерпретации входящих
в них символов) и исследуем, какого рода свойства можно выражать
с помощью формул и какие нельзя.

\section{Формулы и интерпретации}
        \label{formulas-and-structures}%
Начнём с примера. Пусть $M$\т некоторое непустое множество, а
$R$\т бинарное отношение на нём, то есть подмножество
декартова произведения $M\hm\times M$. Вместо $\langle x,y\rangle\hm\in R$
мы будем писать $R(x,y)$.
Рассмотрим формулу
        $$
\forall x\, \exists y\, R(x,y).
        $$
Эта формула выражает некоторое свойство бинарного отношения $R$
(для любого элемента $x\in M$ найдётся элемент, находящийся с
ним в отношении $R$) и может быть истинна или ложна. Например,
если $M$ есть множество натуральных чисел $\bbN$, а $R$\т
отношение \лк строго меньше\пк\ (другими словами, $R$ есть
множество всех пар $\langle x,y\rangle$, для которых $x\hm<y$),
то эта формула истинна. А для отношения \лк строго больше\пк\
(на том же множестве) эта формула ложна.

Вопрос о том, будет ли истинна формула
        $$
\exists y\, R(x,y)
        $$
для данного множества $M$ и для данного бинарного отношения $R$
на нём, не имеет смысла, пока не уточнено, каково значение
переменной $x$. Например, если $M=\bbN$ и $R(x,y)$ есть
$x\hm>y$, то эта формула будет истинной при~$x=3$ и ложной при
$x=0$. Для данных $M$ и $R$ она задаёт некоторое свойство
элемента $x$ и тем самым определяет некоторое подмножество
множества~$M$.

Перейдём к формальным определениям. Пусть $M$\т непустое
множество. Множество~$M^k$ состоит из всех кортежей (последовательностей)
$\langle m_1,\dots,m_k\rangle$ длины~$k$, составленных из
элементов множества~$M$. Назовём \emph{$k$\д местной функцией}%
\index{$k$\д местная функция|defin}%
\index{Функция!$k$\д местная|defin}
на множестве $M$ любое отображение $M^k$ в $M$ (определённое на
всём $M^k$). Синонимы: \лк функция $k$ аргументов\пк, \лк
функция валентности $k$\пк, \лк функция местности $k$\пк\ и даже
\лк функция арности $k$\пк\ (последнее слово происходит от слов
\лк унарная\пк\ для функций одного аргумента, \лк бинарная\пк\
(операция) для функций двух аргументов и \лк тернарная\пк\ для
трёх аргументов).

Назовём \emph{$k$\д местным предикатом}%
\index{$k$\д местный предикат|defin}%
\index{Предикат!$k$\д местный|defin}
на множестве $M$ любое
отображение $M^k$ в множество $\bbB=\{\И,\Л\}$. Такой предикат
будет истинным на некоторых наборах $\langle
m_1,\dots,m_k\rangle$ множества~$M$ и ложным на остальных
наборах. Поставив ему в соответствие множество тех наборов, где
он истинен, мы получаем взаимно однозначное соответствие между
$k$\д местными предикатами на~$M$ и подмножествами множества
$M^k$. Говоря о предикатах, также употребляют термины \лк
валентность\пк, \лк число аргументов\пк\ и др.

Мы будем
рассматривать также функции и предикаты валентности нуль.
Множество $M^0$ одноэлементно (содержит единственную
последовательность длины $0$). Поэтому функции $M^0\to M$
отождествляются с элементами множества $M$, а нульместных
предикатов ровно два\т истинный и ложный.

Естественно, что в формулы будут входить не сами функции и
предикаты, а обозначения для них, которые называют
\emph{функциональными}%
\index{Функциональный символ|defin}%
\index{Символ!функциональный|defin}
и \emph{предикатными символами}.%
\index{Предикатный символ|defin}%
\index{Символ!предикатный|defin}
В качестве символов можно использовать любые знаки. Важно лишь,
что каждому символу приписана валентность, которая определяет,
со сколькими аргументами он может встречаться в формуле.
Произвольный набор предикатных и функциональных символов, для
каждого из которых указано неотрицательное число, называемое
\emph{валентностью},%
\index{Валентность символа|defin}%
\index{Символ!валентность|defin}
мы будем называть \emph{сигнатурой}.%
\index{Сигнатура|defin}

Остаётся определить три вещи: что такое формула данной
сигнатуры, что такое интерпретация данной сигнатуры и когда
формула является истинной (в данной интерпретации).

Фиксируем некоторый набор символов, называемых \emph{индивидными
переменными}.%
\index{Индивидная переменная|defin}%
\index{Переменная!индивидная|defin}
Они предназначены для обозначения элементов
множества, на котором определены функции и предикаты; обычно
в таком качестве используют латинские буквы
$x,y,z,u,v,w$ с индексами. В каждой формуле будет использоваться
конечное число переменных, так что счётного набора переменных
нам хватит. Мы предполагаем, что переменные отличны от всех
функциональных и предикатных символов сигнатуры (иначе выйдет
путаница).

Определим понятие \emph{терма}%
\index{Терм|defin}
данной сигнатуры. Термом называется
последовательность переменных, запятых, скобок и символов
сигнатуры, которую можно построить по следующим правилам:
        %
\begin{itemize}
        \item
Индивидная переменная есть терм.
        \item
Функциональный символ валентности $0$ есть терм.
        \item
Если $t_1,\dots,t_k$\т термы, а $f$\т функциональный символ
валентности $k>0$, то $f(t_1,\dots,t_k)$ есть терм.
        %
\end{itemize}

В принципе можно было не выделять функциональные символы
валентности $0$ (которые также называют \emph{константами})%
\index{Константа|defin}
в отдельную группу, но тогда бы после них
пришлось писать скобки (как это делается в программах на
языке~Си).

Если $A$\т предикатный символ валентности~$k$, а
$t_1,\dots,t_k$\т термы, то выражение $A(t_1,\dots,t_k)$
считается \emph{атомарной формулой}.%
\index{Атомарная формула|defin}%
\index{Формула!атомарная|defin}
Кроме того, любой
предикатный символ валентности~$0$ считается атомарной формулой.

Формулы строятся по таким правилам:
\begin{itemize}
        \item
Атомарная формула есть формула.
        \item
Если $\varphi$\т формула, то $\lnot\varphi$\т формула.
        \item
Если $\varphi$ и $\psi$\т формулы, то выражения $(\varphi
\land\psi)$, $(\varphi\hm\lor\psi)$, $(\varphi\to\psi)$ также
являются формулами.
        \item
Если $\varphi$ есть формула, а $\xi$\т индивидная переменная, то
выражения $\forall \xi \,\varphi$ и $\exists \xi\,\varphi$
являются формулами.
        %
\end{itemize}

Во многих случаях в сигнатуру входит двуместный предикатный
символ $=$, называемый \emph{равенством}.%
\index{Равенство (предикатный символ)|defin}
По традиции вместо $=(t_1,t_2)$ пишут $(t_1=t_2)$.

Итак, понятие формулы в данной сигнатуре полностью определено.
Иногда такие формулы называют \emph{формулами первого порядка}%
\index{Формула!первого порядка|defin}
данной сигнатуры, или формулами \emph{языка первого порядка}%
\index{Языки!первого порядка|defin}
с данной сигнатурой.

Наш следующий шаг\т определение \emph{интерпретации}%
\index{Интерпретация|defin}
данной
сигнатуры. Пусть фиксирована некоторая сигнатура $\sigma$. Чтобы
задать интерпретацию сигнатуры $\sigma$, необходимо:
        %
\begin{itemize}
        \item
указать некоторое непустое множество $M$, называемое
\emph{носителем}%
\index{Интерпретация!носитель|defin}%
\index{Носитель интерпретации|defin}
интерпретации;
        \item
для каждого предикатного символа сигнатуры~$\sigma$ указать
предикат с соответствующим числом аргументов, определённый
на множестве~$M$ (как мы уже говорили,
$0$\д местным предикатным символам ставится в соответствие
либо \И, либо \Л);
        \item
для каждого функционального символа сигнатуры~$\sigma$
указать функцию соответствующего числа аргументов с аргументами
и значениями из~$M$ (в частности, для $0$\д местных функциональных
символов надо указать элемент множества $M$, с ними сопоставляемый).
        %
\end{itemize}

Если сигнатура включает в себя символ равенства, то среди её
интерпретаций выделяют \emph{нормальные}%
\index{Нормальная интерпретация|defin}%
\index{Интерпретация!нормальная|defin}
интерпретации, в
которых символ равенства интерпретируется как совпадение
элементов.


Приведём несколько примеров сигнатур, используемых в различных
теориях.

Сигнатура теории упорядоченных множеств включает
в себя два двуместных предикатных символа (равенство и
порядок) и не имеет функциональных символов. Здесь также
вместо $\le(x,y)$ по традиции пишут $x \le y$.

Аксиомы порядка%
\index{Аксиомы!порядка}
(рефлексивность, антисимметричность,
транзитивность) могут быть записаны формулами этой сигнатуры.
Например, требование антисимметричности записывается так:
        $$
\forall x\,\forall y (((x\le y)\land (y\le x))\to(x=y)).
        $$

Иногда в сигнатуру теории упорядоченных множеств вместо символа
$\le$ включают символ $<$; большой разницы~тут нет.

\begin{problem}
        %
Как записать с помощью формулы свойство линейной
упорядоченности? свойство не иметь наибольшего элемента?
свойство плотности (отсутствия соседних элементов)? свойство
фундированности (отсутствия бесконечных убывающих последовательностей\т
или, что эквивалентно, наличия минимального элемента в любом подмножестве)?
свойство полной упорядоченности? (Указание: не
для всех перечисленных свойств это возможно.)
        %
\end{problem}

Сигнатуру теории групп можно выбирать по\д разному. Можно
считать, что (помимо равенства) она имеет двуместный
функциональный символ $\times$ (который по традиции записывают
между множителями), константу (нульместный функциональный
символ) $1$ и одноместный функциональный символ $\inv(x)$ для
обращения. Тогда аксиомы теории групп записываются с
использованием лишь кванторов всеобщности:
        \begin{align*}
&\forall x\,\forall y\,\forall z
      (((x\times y)\times z)=(x\times(y\times z))),\\
&\forall x\, (((x\times 1)= x)\land((1\times x)=x)),\\
&\forall x\, (((x\times \inv(x))= 1)\land((\inv(x)\times x)=1)).
        \end{align*}

Если не включать операцию обращения в сигнатуру, придётся
использовать квантор существования и переписать последнюю
аксиому так:
        $$
\forall x\, \exists y\,((( x\times y)= 1)\land((y\times x)=1)).
        $$

\begin{problem}
        %
Как записать аксиомы теории групп, если в сигнатуре нет константы~$1$?
(Указание: аксиома о существовании обратного станет частью аксиомы
о существовании единицы.)
        %
\end{problem}

\begin{problem}
        %
Как записать в виде формулы требование коммутативности
группы? утверждение о том, что любой элемент (кроме единицы)
имеет порядок~$11$? конечность группы? (Указание: не всё из
перечисленного можно записать, хотя пока у нас нет средств это
установить.)
        %
\end{problem}

Сигнатура теории множеств содержит два двуместных предикатных
символа: для принадлежности и для равенства. Аксиомы теории
множеств можно записывать в виде формул этой сигнатуры. Чаще
всего рассматривают вариант аксиоматической теории множеств,
называемый теорией Цермело\ч Френкеля%
\index{Теория!Цермело\ч Френкеля}%
\glossary{\Цермело}%
\glossary{\Френкель}
и обозначаемый ZF.%
\index{ZF|defin}
Приведём для примера одну из аксиом теории ZF, называемую
\emph{аксиомой объёмности}, или \emph{экстенсиональности}:%
\index{Аксиома!объёмности (экстенсиональности)|defin}
$$
\forall x\,\forall y\,((\forall z\,((z\in x)\to(z\in y))\land
                        \forall z\,((z\in y)\to(z\in x))) \to (x=y)).
$$

\begin{problem}
        %
Сформулировать словесно эту аксиому.
        %
\end{problem}

\begin{problem}
        %
Записать в виде формулы \emph{аксиому регулярности}, или \emph{фундирования},%
\index{Аксиома!регулярности (фундирования)|defin}
которая говорит, что у всякого множества есть минимальный (с точки зрения
отношения $\in$) элемент, то есть элемент, не пересекающийся с самим
множеством.
        %
\end{problem}

\begin{problem}
        %
Какова естественная сигнатура для теории полей? Можно ли
записать в виде формулы этой сигнатуры утверждение о том, что
поле имеет характеристику~$2$? конечную характеристику?
алгебраически замкнуто?
        %
\end{problem}

\section{Определение истинности}
        \label{truth-definition}%
Из приведённых выше примеров, вероятно, понятен смысл формулы,
то есть ясно, в каких интерпретациях данной сигнатуры и для
каких элементов формула истинна. Тем не менее для любителей
строгости мы приведём формальное определение истинности. (Его
детали понадобятся, когда мы будем проверять истинность
выводимых формул, см.~раздел~\ref{predicate-correctness}.)

Прежде всего, определим формально понятие \emph{параметра}%
\index{Параметр формулы|defin}%
\index{Формула!параметр|defin}
формулы (переменной, от значения которой может зависеть
истинность формулы). Согласно этому определению формула
        $
\forall x\,\exists y\, A(x,y)
        $
не имеет параметров, а формулы
        $
 \exists y\, A(x,y)
        $
и
        $
(A(x)\land \forall x\, B(x,x))
        $
имеют единственный параметр $x$. Вот как выглядит это
определение:
        %
\begin{itemize}
        \item
Параметрами терма являются все входящие в него
индивидные переменные.
        \item
Параметрами атомарной формулы являются параметры всех
входящих в неё термов.
        \item
Параметры формулы $\lnot \varphi$ те же, что у формулы
$\varphi$.
        \item
Параметрами формул $(\varphi\land\psi)$, $(\varphi\lor\psi)$ и
$(\varphi\to\psi)$ являются все параметры формулы $\varphi$, а
также все параметры формулы $\psi$.
        \item
Параметрами формул $\forall\xi\,\varphi$ и $\exists\xi\,\varphi$
являются все параметры формулы $\varphi$, кроме переменной $\xi$.
        %
\end{itemize}

Параметры иногда называют \emph{свободными переменными}%
\index{Свободная переменная|defin}%
\index{Переменная!свободная|defin}
формулы.
Заметим, что формула может иметь одновременно параметр $x$ и
использовать (в другом месте) квантор $\forall x$. Как говорят в
этом случае, одна и та же переменная имеет \emph{свободные}%
\index{Свободное вхождение переменной|defin}%
\index{Переменная!свободное вхождение|defin}
и
\emph{связанные}%
\index{Связанное вхождение переменной|defin}%
\index{Переменная!связанное вхождение|defin}
вхождения. Свободное вхождение переменной\т это
такое вхождение, которое не входит в область действия одноимённого
квантора. Если аккуратно определить эту область действия,
несложно проверить, что параметры формулы\т это как раз
переменные, имеющие свободные вхождения.

Теперь мы хотим определить понятие формулы, истинной в данной
интерпретации при данных значениях параметров. Технически проще
считать, что всем индивидным переменным приписаны какие\д то
значения, а потом доказать, что переменные, не являющиеся
параметрами, не влияют на истинность формулы.

Итак, пусть фиксирована сигнатура и некоторая интерпретация
этой сигнатуры. \emph{Оценкой}%
\index{Оценка|defin}
назовём отображение, которое
ставит в соответствие каждой индивидной
переменной некоторый элемент
множества, являющегося носителем интерпретации. Этот элемент
будем называть \emph{значением переменной}%
\index{Значение переменной|defin}%
\index{Переменная!значение|defin}
при данной оценке.

Определим индуктивно \emph{значение терма}%
\index{Значение терма|defin}%
\index{Терм!значение|defin}
$t$ при данной оценке
$\pi$, которое мы будем обозначать $[t](\pi)$.
        \begin{itemize}\item
Для переменных оно уже определено.
        \item
Если $t$ является константой (нульместным
функциональным символом), то $[t](\pi)$ не зависит от $\pi$ и
равно значению этой константы при данной интерпретации
(напомним, в интерпретации с каждой константой сопоставляется
некоторый элемент но\-си\-теля).
        \item
Если $t$ имеет вид $f(t_1,\dots,t_m)$, где $f$\т функциональный
символ валентности $m$, а $t_1,\dots,t_m$\т термы, то
$[t](\pi)$ определяется как $[f]([t_1](\pi),\dots,[t_m](\pi))$, где
$[f]$ есть функция, соответствующая символу $f$ в нашей
интерпретации, а $[t_i](\pi)$ есть значение терма $t_i$
при оценке $\pi$.
        \end{itemize}

Теперь можно определить \emph{значение формулы}%
\index{Значение формулы|defin}%
\index{Формула!значение|defin}
$\varphi$ при
данной оценке $\pi$ в данной интерпретации, которое обозначается
$[\varphi](\pi)$ и может быть равно
\И\ или \Л; в первом случае
формула называется \emph{истинной},%
\index{Истинная формула|defin}%
\index{Формула!истинная|defin}
во втором\т \emph{ложной}.%
\index{Ложная формула|defin}%
\index{Формула!ложная|defin}
Это определение также индуктивно:
       \begin{itemize}
       \item
Значение атомарной формулы $A(t_1,\dots,t_m)$ определяется
как $[A]([t_1](\pi),\dots,[t_m](\pi))$, где $[A]$\т
предикат, соответствующий предикатному символу $A$ в
рассматриваемой интерпретации. Если формула
представляет собой нульместный предикатный символ, то её
значение не зависит от оценки и есть значение этого
символа.
        \item
$[\lnot \varphi](\pi)$ определяется как $\lnot [\varphi(\pi)]$,
где $\lnot$ понимается как операция в $\bbB$. Другими словами,
формула $\lnot \varphi$ истинна при оценке $\pi$ тогда и только
тогда, когда формула $\varphi$ ложна при этой оценке.
        \item
$[\varphi \land \psi](\pi)$ определяется как
$[\varphi](\pi)\land [\psi](\pi)$, где $\land$ в правой
части понимается как операция в $\bbB$. (Другими словами,
формула $(\varphi\land\psi)$ истинна при оценке $\pi$
тогда и только тогда, когда обе формулы $\varphi$ и
$\psi$ истинны при этой оценке.) Аналогичным образом
$[\varphi \lor \psi](\pi)$ определяется как
$[\varphi](\pi)\lor [\psi](\pi)$,  а
$[\varphi \to \psi](\pi)$\т как
$[\varphi](\pi)\to [\psi](\pi)$.
        \item
Формула $\forall \xi\,\varphi$ истинна на оценке $\pi$ тогда и
только тогда, когда формула $\varphi$ истинна на любой оценке
$\pi'$, которая совпадает с $\pi$ всюду, кроме значения
переменной $\xi$ (которое в оценке~$\pi'$
может быть любым).
Другими словами, если обозначить через $\pi+(\xi\mapsto m)$
оценку, при которой значение переменной $\xi$ равно $m$,
а остальные переменные принимают те же значения, что и
в оценке $\pi$, то
        $$
[\forall \xi\, \varphi](\pi)=
      \bigwedge_{m\in M} [\varphi](\pi+(\xi\mapsto m)).
        $$
(В правой части стоит бесконечная конъюнкция,
которая истинна, если все её члены истинны.)
        \item
Формула $\exists \xi\,\varphi$ истинна на оценке $\pi$ тогда и
только тогда, когда формула $\varphi$ истинна на некоторой оценке
$\pi'$, которая совпадает с $\pi$ всюду, кроме значения
переменной $\xi$ (которое в оценке~$\pi'$
может быть любым). Другими словами,
        $$
[\exists \xi\, \varphi](\pi)=
      \bigvee_{m\in M} [\varphi](\pi+(\xi\mapsto m)).
        $$
(В правой части стоит бесконечная дизъюнкция,
которая истинна, если хотя бы один из её членов
истинен.)
        %
\end{itemize}

Заметим, что в двух последних пунктах значение переменной $\xi$
в оценке $\pi$ не играет роли. Это позволяет легко доказать
(индукцией по построению формулы) такое утверждение: если две
оценки $\pi_1$ и $\pi_2$ придают одинаковые значения всем
параметрам формулы $\varphi$, то $[\varphi](\pi_1)\hm=
[\varphi](\pi_2)$. Другими словами, истинность формулы
определяется значениями её параметров.

        \begin{problem}
Проведите это индуктивное рассуждение подробно.
        \end{problem}

        \begin{problem}
Приведённые выше определения применимы к любой формуле, в том
числе и к странной формуле $\forall y\, A(x)$. Какие у неё
параметры? При каких значениях параметров она истинна?
(Ответ: она имеет единственный параметр~$x$ и эквивалентна
формуле~$A(x)$.)
        \end{problem}

        \begin{problem}
В каком случае будет истинна формула $\forall x\,\exists
x\,A(x)$? Тот же вопрос для формулы $\exists x\,\forall x\,
A(x)$. (Ответ: первая из этих формул эквивалентна формуле
$\exists x\, A(x)$, а вторая\т формуле~$\forall x\, A(x)$.)
        \end{problem}

Формула называется \emph{замкнутой},%
\index{Замкнутая формула|defin}%
\index{Формула!замкнутая|defin}
если она не имеет
параметров. Замкнутые формулы называют также \emph{суждениями}.%
\index{Суждение|defin}
Как мы доказали, истинность замкнутой формулы определяется
выбором интерпретации (и не зависит от значений переменных).

\section{Выразимые предикаты}
        \label{definable-predicates}%
Пусть фиксирована некоторая сигнатура $\sigma$ и её
интерпретация с носителем $M$. Мы хотим определить понятие
\emph{выразимого}
(с помощью формулы данной сигнатуры в данной
интерпретации) $k$\д местного предиката.

Выберем $k$ переменных $x_1,\dots, x_k$ и рассмотрим произвольную
формулу $\varphi$, все параметры которой содержатся в списке
$x_1,\dots,x_k$. Истинность этой формулы зависит только от
значений переменных $x_1,\dots,x_k$. Тем самым возникает
отображение $M^k\to\bbB=\{\И,\Л\}$, то есть некоторый $k$\д
местный предикат на $M$. Говорят, что этот предикат
\emph{выражается}%
\index{Предикат!выражаемый формулой|defin}
формулой $\varphi$. Все предикаты, которые
можно получить таким способом, называются \emph{выразимыми}.%
\index{Выразимый предикат|defin}%
\index{Предикат!выразимый|defin}
(Ясно, что конкретный выбор списка переменных роли не играет.)
Соответствующие им подмножества множества $M^k$ (области
истинности выразимых предикатов) также называют выразимыми.

\begin{problem}
        %
Докажите, что пересечение, объединение и разность двух выразимых
множеств являются выразимыми. Докажите, что проекция $k$\д
мерного выразимого множества вдоль одной из \лк осей
координат\пк\ является $(k-1)$\д мерным выразимым множеством.
        %
\end{problem}


\textsf{Пример}.  Рассмотрим сигнатуру, содержащую одноместный
функциональный символ $S$ и двуместный предикатный символ
равенства ($=$), и интерпретацию этой сигнатуры. В
качестве носителя возьмём натуральный ряд $\bbN$. Символ~$S$
будет обозначать функцию увеличения на~$1$ (можно считать $S$
сокращением от слова successor\т последователь).  Равенство
интерпретируется как совпадение элементов.

Легко проверить, что одноместный предикат \лк быть нулём\пк\
выразим в этой интерпретации, несмотря на то, что константы для
нуля в сигнатуре не предусмотрено. В самом деле, он выражается
формулой
        $$
\lnot\exists y (x=S(y))
        $$
с единственным параметром $x$.

Ещё проще выразить двуместный предикат \лк быть
больше на $2$\пк, при этом даже не нужны кванторы: $y=S(S(x))$.

Любопытно, что уже в такой простой ситуации можно сформулировать
содержательную задачу: выразить предикат $y=x+N$, где $N$\т
большое число (скажем, миллиард), с помощью существенно более
короткой формулы, чем $y=S(S(\dots(S(x))\dots))$. Как ни
удивительно, это вполне возможно, и соответствующую формулу
вполне можно уместить на листе бумаги.
        %
\begin{problem}
        %
Докажите, что предикат $y=x+N$ можно выразить формулой указанной
сигнатуры, длина которой есть $O(\log N)$. (Указание. Если мы
научились выражать $y=x+n$, можно выразить $y=x+2n$ с помощью
формулы
        $$
\exists z\, ((z=x+n)\land(y=z+n))
        $$
(в которой через $z=x+n$ и $y=z+n$ обозначены соответствующие
формулы). Это само по себе ничего не даёт, так как длина
формулы увеличилась вдвое, но можно использовать такой трюк:
        $$
\exists z\, \forall u\,\forall v
(((u=x\land v=z)\lor(u=z\land v=y))\to(v=u+n)).
        $$
Далее можно воспользоваться записью числа~$N$ в двоичной системе
счисления.)
        %
\end{problem}

Можно доказать, что в этой сигнатуре кванторы почти не увеличивают
набор выразимых предикатов: всякий выразимый предикат будет выражаться
бескванторной формулой (возможно, гораздо более длинной),
если добавить к сигнатуре константу~$0$.
Мы вернёмся к этому вопросу в разделе~\ref{quantifier-elimination}.

Чтобы привыкнуть к понятию выразимости, рассмотрим ещё и такой
пример. Пусть сигнатура содержит предикат равенства и
трёхместный предикат $C$. Рассмотрим интерпретацию, в которой
носителем является множество точек плоскости, равенство
интерпретируется как совпадение точек, а $C(x,y,z)$ означает,
что точки $x$ и $y$ равноудалены от точки $z$. Оказывается,
что этого предиката достаточно, чтобы выразить более или менее
все традиционные понятия элементарной геометрии.

Как, например, записать, что три различные точки $A,B,C$ лежат
на одной прямой? Вот как: \лк не существует другой точки $C'$,
которая находилась бы на тех же расстояниях от $A$ и $B$, что и
точка $C$\пк.
        %
\begin{problem}
        %
Напишите соответствующую формулу указанной сигнатуры.
        %
\end{problem}

Теперь легко выразить такое свойство четырёх точек $A,B,C,D$:
\лк точки $A$ и $B$ различны, точки $C$ и $D$ различны и
прямые $AB$ и $CD$ параллельны\пк. В самом деле, надо написать,
что нет точки, которая бы одновременно лежала на одной прямой
с $A$ и $B$, а также на одной прямой с $C$ и $D$.

После этого можно выразить свойство четырёх точек \лк быть вершинами
параллелограмма\пк. Это позволяет переносить отрезок
параллельно себе. После этого легко выразить и такое свойство:
\лк расстояние $AB$ равно расстоянию $CD$\пк.

\begin{problem}
        %
Запишите соответствующую формулу.
        %
\end{problem}

\smallskip
Аналогичным образом можно двигаться и дальше.

\smallskip
\begin{problem}
        %
Выразите свойство $|OA|\le |OB|$ трёх точек $O$, $A$,~$B$.
(Указание. Напишите, что все прямые, проходящие через~$A$,
пересекаются с окружностью радиуса~$OB$ с центром в~$O$.)
        %
\end{problem}

\begin{problem}
        %
Запишите в виде формулы:
(\textbf{а})~равенство треугольников;
(\textbf{б})~равенство углов;
(\textbf{в})~свойство угла быть прямым.
        %
\end{problem}

\begin{problem}
        %
Рассмотрим естественную интерпретацию сигнатуры $({=},{<})$ на
множестве целых чисел. Как выразить предикат $y=x+1$?
        %
\end{problem}

\begin{problem}
        %
Рассмотрим множество действительных чисел как интерпретацию
сигнатуры $({=},{+}, {y=x^2})$. Как выразить трёхместный
предикат $xy=z$?
        %
\end{problem}

\begin{problem}
        %
Рассмотрим множество целых положительных чисел как интерпретацию
сигнатуры, содержащей равенство и двуместный предикат \лк $x$
делит $y$\пк. Выразите свойства \лк равняться единице\пк\ и \лк
быть простым числом\пк.
        %
\end{problem}

\begin{problem}
        %
Рассмотрим плоскость как интерпретацию сигнатуры, содержащей
предикат равенства (совпадение точек) и двуместный предикат \лк
находиться на расстоянии~$1$\пк. Выразите двуместные предикаты
\лк находиться на расстоянии~$2$\пк\ и \лк находиться на
расстоянии не более~$2$\пк. Выразите двуместный предикат \лк
находиться на расстоянии $1/2$\пк.
        %
\end{problem}

\section{Выразимость в арифметике}
        \label{arithmetic-definability}%

Рассмотрим сигнатуру, имеющую два двуместных функциональных
символа\т сложение и умножение (как обычно, мы будем писать
$x+y$ вместо $+(x,y)$ \итд) и двуместный предикатный символ
равенства. Рассмотрим интерпретацию этой сигнатуры, носителем
которой является множество $\bbN$ натуральных чисел, а сложение,
умножение и равенство интерпретируются стандартным образом.

Выразимые с помощью формул этой сигнатуры предикаты называются
\emph{арифметическими}%
\index{Арифметический предикат|defin}%
\index{Предикат!арифметический|defin}
и играют в математической логике важную
роль. Соответствующие множества также называются
\emph{арифметическими}.%
\index{Арифметическое множество|defin}%
\index{Множество!арифметическое|defin}
О них подробно рассказано в другой нашей
книжке~\cite{computable-functions}; оказывается, что почти
всякое множество, которое можно описать словами, является
арифметическим.

\begin{problem}
        %
Докажите, что существует множество натуральных чисел, не
являющееся арифметическим. (Указание: семейство всех подмножеств
множества~$\bbN$ несчётно, а арифметических множеств
счётное число.)
        %
\end{problem}

Для начала мы установим арифметичность довольно простых
предикатов.

\begin{itemize}
        \item
Предикат $x\le y$ является арифметическим. В самом деле, его
можно записать как $\exists z\, (x+z=y)$.
        \item
Предикаты $x=0$ и $x=1$ являются арифметическими. В самом деле,
$x=0$ тогда и только тогда, когда $x\le y$ для любого $y$ (а
также когда $x+x=x$). А~$x=1$ тогда и только
тогда, когда $x$ представляет собой наименьшее число, отличное
от нуля. (Можно также воспользоваться тем, что $y\cdot 1=y$ при
любом~$y$.)
        \item
Вообще для любого фиксированного
числа~$c$ предикат ${x=c}$ является арифметическим.
(Например, можно написать сумму из большого числа единиц.)
        \item
Полезно такое общее наблюдение: если мы уже установили, что
какой\д то предикат является арифметическим, то в дальнейшем его
можно использовать в формулах, как если бы он входил в
сигнатуру, поскольку его всегда можно заменить на выражающую его
формулу.
        \item
Предикат $x|y$ (число $x$ является делителем числа $y$), очевидно,
арифметичен (формула $\exists z\, (xz=y)$).
        \item
Предикат \лк $x$\т простое число\пк\ арифметичен. В самом деле, число
просто, если оно отлично от~$1$ и
любой его делитель равен $1$ или самому числу. Это
сразу же записывается в виде формулы.
        \item
Операции частного и остатка арифметичны (в том смысле, что
трёхместные предикаты \лк $q$ есть частное при делении $a$ на $b$\пк\
и \лк $r$ есть остаток при делении $a$ на $b$\пк\ арифметичны.
Например, первый из них записывается формулой
        $
\exists r\, ((a=bq+r)\land (r<b))
        $
(как мы уже говорили, использование арифметического предиката
$(r<b)$ не создаёт проблем).
        \item
Этот список можно продолжать: для многих предикатов их определение
по существу уже является нужной формулой. Например, свойства
\лк быть наибольшим общим делителем\пк, \лк быть наименьшим
общим кратным\пк, \лк быть взаимно простыми\пк\ все относятся
к этой категории.
        \item
Предикат \лк быть степенью двойки\пк\ является арифметическим
(хотя это и не столь очевидно, как в предыдущих примерах).
В самом деле, это свойство можно переформулировать так: любой
делитель либо равен единице, либо чётен.
        %
\end{itemize}

Последнее из наших рассуждений годится для степеней тройки и
вообще для степеней любого простого числа. Однако, скажем, для
степеней шестёрки оно не проходит, и, пожалуй, мы подошли к
границе, где без некоторого общего метода не обойтись.

Два наиболее известных способа доказывать арифметичность
основаны на возможности \лк кодирования\пк\ конечных
множеств и последовательностей. Один из них восходит к Гёделю%
\glossary{\Гёдель}
(так называемая $\beta$\д функция Гёделя),%
\index{$\beta$\д функция Гёделя}
второй изложен в
книге \лк Теория формальных систем\пк~\cite{smullyan-formal}. Её
написал Р.\,Смаллиан,%
        \glossary{\Смаллиан}
известный также как автор популярных
сборников \лк логических задач\пк\ и анекдотов. (Один из таких
сборников имеет парадоксальное название \лк Как же называется
эта книга?\пк~\cite{smullyan-informal}.)

В некоторых отношениях метод Гёделя предпочтительней, и мы
рассказываем о нём в книжке о вычислимых
функциях~\cite{computable-functions}, но сейчас для разнообразия
рассмотрим другой способ.%
\index{Кодирование последовательностей и множеств}
Зафиксируем взаимно однозначное
соответствие между натуральными числами и двоичными словами:
        $$
\begin{array}{cccccccccc}
0 & 1 & 2 & 3 & 4 & 5 & 6 & 7 & 8&\dots \\
\Lambda& 0 & 1 & 00 & 01 & 10 & 11 & 000 & 001 &\ldots
\end{array}
        $$
Это соответствие задаётся так: чтобы получить слово,
соответствующее числу $n$, надо записать $n+1$ в двоичной
системе и удалить первую единицу. Например, нулю соответствует
пустое слово~$\Lambda$, числу~$15$\т слово $0000$ \итд Теперь
можно говорить об арифметичности предикатов, определённых на
двоичных словах, имея в виду арифметичность соответствующих
предикатов на $\bbN$.

\begin{itemize}
      \item
Предикат \лк слово $x$ состоит из одних нулей\пк\ арифметичен. В
самом деле, при переходе к числам ему соответствует предикат
\лк$x+1$ есть степень двойки\пк, который (как мы видели)
арифметичен.
        \item
Предикат \лк слова $x$ и $y$ имеют одинаковую длину\пк\
арифметичен. В самом деле, это означает, что найдётся степень
двойки $c$, для которой $c-1\hm\le x,y\hm<2c-1$ (именно такой
промежуток заполняют числа, которым соответствуют слова одной
длины). Аналогичным образом устанавливается арифметичность
предиката \лк $x$ короче $y$\пк.
        \item
Предикат \лк слово $z$ является конкатенацией слов $x$ и $y$\пк\
(проще говоря, $z$ получается приписыванием $y$ справа к слову $x$)
арифметичен. В самом деле, его можно выразить так: найдётся
слово $y'$ из одних нулей, имеющее ту же длину, что и слово $y$,
при этом $(z+1)\hm=(x+1)(y'+1)+(y-y')$ (умножение на $y'+1$
соответствует дописыванию нулей, а добавление $y-y'$ заменяет
нули на буквы слова $y$).
        \item
Предикат \лк слово $x$ является началом слова $y$\пк\
арифметичен. В самом деле, это означает, что существует слово
$t$, при котором $y$ есть конкатенация $x$ и $t$.
        \item
То же самое верно для предикатов \лк $x$ есть конец слова $y$\пк,
\лк $x$ есть подслово слова $y$\пк\ (последнее означает, что
найдутся слова $u$ и $v$, для которых $y$ есть конкатенация
$u$, $x$ и $v$; конкатенация трёх слов выразима через
конкатенацию двух).
        \item
Существует арифметический трёхместный предикат $S(x,a,b)$, обладающий
такими свойствами: (\textbf{а})~для любых $a$ и $b$ множество
$S_{ab}=\{x\mid S(x,a,b)\}$%
\index{$S_{ab}$|defin}
конечно; (\textbf{б})~среди множеств
$S_{ab}$ при различных парах $a,b$ встречаются все конечные
множества. Например, в качестве такого предиката можно взять \лк
$x$ короче $a$ и
$axa$ есть подслово слова $b$\пк\ (здесь $axa$ есть конкатенация
трёх слов: $a$, $x$ и снова $a$).

В самом деле, ясно, что слово $x$ не длиннее слова $b$, и потому
множество $S_{ab}$ всегда конечно. С~другой стороны, пусть
имеется некоторое конечное множество слов $x_1,\dots, x_n$.
Положим $a=100\dots001$, где число нулей больше длины любого из
слов $x_i$, и $b\hm=ax_1ax_2a\dots ax_na$.
        %
\end{itemize}

Последнее утверждение не упоминает слова, и больше они
нам не понадобятся: достаточно знать, что конечные множества
натуральных чисел можно кодировать парами натуральных чисел
в описанном смысле.

Теперь мы можем выразить, что число $x$ является степенью числа
$4$, следующим образом: существует конечное множество $U$,
которое содержит число $x$ и обладает таким свойством: всякий
элемент $u\in U$ либо равен $1$, либо делится на $4$ и $u/4$
также принадлежит $U$. Теперь надо везде заменить множество $U$
на его код $u_1,u_2$, а утверждение $x\in U$ на $S(x,u_1,u_2)$,
где $S$\т построенный нами кодирующий предикат.

Немного сложнее выразить двуместный предикат $x\hm=4^k$. Тут
хотелось бы сказать так: существует последовательность
$x_0,x_1,\dots,x_k$, для которой $x_0=1$, каждый следующий член
вчетверо больше пре\-дыдущего ($x_{i+1}\hm=4x_i$) и $x_k\hm= x$.
Как научиться говорить о последовательностях, если мы умеем
говорить о множествах? Вспомним, что в терминах теории множеств
последовательность есть функция, определённая на начальном
отрезке натурального ряда, то есть конечное множество пар
$\{\langle 0,x_0\rangle, \langle 1, x_1\rangle,\dots, \langle k,
x_k\rangle\}$. Пары можно кодировать числами. Например, можно
считать кодом пары $\langle x,y\rangle$ число
$c\hm=(x+y)^2+x$, поскольку по нему арифметически
восстанавливается $x+y$ (как наибольшее число, квадрат которого
не превосходит $c$), а затем~$x$ и~$y$. Теперь конечное
множество пар можно заменить конечным множеством их кодов,
которое в свою очередь можно закодировать парой чисел.

\begin{problem}
        %
Проведите это рассуждение подробно.
        %
\end{problem}

\begin{problem}
        %
Покажите, что двуместный предикат \лк $x$ есть $n$\д ое
по порядку простое число\пк\ арифметичен.
        %
\end{problem}

\section{Невыразимые предикаты: автоморфизмы}
        \label{automorphisms}
Мы видели, как можно доказать выразимость некоторых свойств.
Сейчас мы покажем, каким образом можно доказывать
невыразимость.

Начнём с такого примера. Пусть
сигнатура содержит двуместный предикат равенства
($=$) и двуместную операцию сложения ($+$). Рассмотрим её
интерпретацию, носителем которой являются целые числа, а
равенство и сложение интерпретируются стандартным образом.
Оказывается, что предикат ${x>y}$ не является выразимым.

Причина очевидна: с точки зрения сложения целые числа устроены
симметрично, положительные ничем не отличаются от отрицательных.
Если мы изменим знак у всех переменных, входящих в формулу, то
её истинность не может измениться. Но при этом $x>y$ заменится
на $x<y$, и потому это свойство не является выразимым.

Формально говоря, надо доказывать по индукции такое свойство:
если формула $\varphi$ указанной сигнатуры
истинна при оценке $\pi$, то она истинна и
при оценке $\pi'$, в которой значения всех переменных меняют
знак. (Подробно мы объясним это в общей ситуации дальше.)

Сформулируем общую схему, которой следует это рассуждение. Пусть
имеется некоторая сигнатура $\sigma$ и интерпретация этой
сигнатуры, носителем которой является множество $M$. Взаимно
однозначное отображение $\alpha\colon M\to M$ называется
\emph{автоморфизмом}%
\index{Автоморфизм интерпретации|defin}%
\index{Интерпретация!автоморфизм|defin}
интерпретации, если все функции и
предикаты, входящие в интерпретацию, устойчивы относительно
$\alpha$. При этом $k$\д местный предикат~$P$ называется
\emph{устойчивым}%
\index{Устойчивый предикат|defin}%
\index{Предикат!устойчивый|defin}
относительно~$\alpha$, если
        $$
P(\alpha(m_1),\dots,\alpha(m_k))
\Leftrightarrow
P(m_1,\dots, m_k)
        $$
для любых элементов
$m_1,\dots,m_k\hm\in M$. Далее, $k$\д местная
функция $f$ называется \emph{устойчивой}%
\index{Устойчивая функция|defin}%
\index{Функция!устойчивая|defin}
относительно $\alpha$, если
        $$
f(\alpha(m_1),\dots,\alpha(m_k))=
\alpha(f(m_1,\dots, m_k)).
        $$
Это определение обобщает стандартное определение
автоморфизма для групп, колец, полей \итд

\begin{theorem}
        \label{nondefinability-authomorphisms}
Любой предикат, выразимый в данной интерпретации,
устойчив относительно её автоморфизмов.
        %
\end{theorem}

\begin{proof}
        %
Проведём доказательство этого (достаточно очевидного) утверждения
формально.

Пусть $\pi$\т некоторая оценка, то есть отображение, ставящее в
соответствие всем индивидным переменным некоторые элементы
носителя. Через $\alpha\circ\pi$ обозначим оценку, которая
получится, если к значению каждой переменной применить
отображение $\alpha$; другими словами,
$\alpha\circ\pi(\xi)=\alpha(\pi(\xi))$ для любой переменной
$\xi$.

Первый шаг состоит в том, чтобы индукцией по построению терма
$t$ доказать такое утверждение: значение терма $t$ при оценке
$\alpha\circ\pi$ получается применением $\alpha$ к значению
терма $t$ при оценке $\pi$:
        $$
[t](\alpha\circ\pi)=\alpha([t](\pi)).
        $$
Для переменных это очевидно, а шаг индукции использует
устойчивость всех функций интерпретации относительно $\alpha$.

Теперь индукцией по построению формулы $\varphi$ легко доказать такое
утверждение:
        $$
[\varphi](\alpha\circ\pi) = [\varphi] (\pi).
        $$
Мы не будем выписывать эту проверку; скажем лишь, что
взаимная однозначность $\alpha$ используется, когда мы разбираем
случай кванторов. (В самом деле, если с одной стороны изоморфизма
берётся какой\д то объект, то взаимная однозначность позволяет
взять соответствующий ему объект с другой стороны изоморфизма.)
        %
\end{proof}

Теорема~\ref{nondefinability-authomorphisms} позволяет доказать
невыразимость какого\д то предиката, предъявив автоморфизм
интерпретации, относительно которого интересующий нас предикат
неустойчив. Вот несколько примеров:

\begin{itemize}
        \item
$(\bbZ,{=},{<})$ %\quad
Сигнатура содержит равенство и отношение порядка. Интерпретация:
целые числа. Невыразимый предикат: $x=0$. Автоморфизм: $x\mapsto
x+1$.
        \item
$(\bbQ,{=},{<},{+})$ %\quad
Сигнатура содержит равенство, отношение порядка и операцию
сложения. Интерпретация: рациональные числа. Невыразимый
предикат: $x=1$. Автоморфизм: $x\hm\mapsto 2x$.

Заметим, что сложение позволяет выразить предикат $x=0$.
Кроме того, отметим, что вместо рациональных чисел можно взять
действительные (но не целые, так как в этом случае единица
описывается как наименьшее число, большее нуля).

        \item
$(\bbR,{=},{<},0,1)$ %\quad
Сигнатура содержит равенство, порядок и константы $0$ и~$1$.
Интерпретация: действительные числа. Невыразимый предикат:
$x=1/2$. (Автоморфизм упорядоченного множества $\bbR$, сохраняющий
$0$ и~$1$, но не $1/2$, построить легко.)
        \item
$(\bbR,{=},{+}, 0, 1)$ %\quad
Сигнатура содержит равенство, сложение, константы $0$ и~$1$.
Интерпретация: действительные числа. Одноместный предикат
$x=\gamma$ выразим для рациональных~$\gamma$ и невыразим для
иррациональных~$\gamma$.

В самом деле, выразимость для рациональных $\gamma$ очевидна.
Не\-вы\-ра\-зи\-мость для иррациональных $\gamma$ следует из того, что
для любых двух иррациональных~$\gamma_1$ и~$\gamma_2$ существует
автоморфизм, переводящий~$\gamma_1$ в~$\gamma_2$. (В самом деле,
рассмотрим $\bbR$ как бесконечномерное векторное пространство
над $\bbQ$. Векторы~$1,\gamma_1$ линейно независимы и потому их
можно дополнить до базиса Гамеля (подробности смотри в книжке по
теории множеств~\cite{naive-set-theory}). Сделаем то же самое с
векторами $1,\gamma_2$. Получатся равномощные базисы, после чего
мы берём $\bbQ$\д линейный оператор, переводящий $1$ в $1$ и
$\gamma_1$ в $\gamma_2$.)
        \item
$(\bbC,{=},{+},{\times},0,1)$ %\quad
В сигнатуру входят предикат ра\-вен\-с\-т\-ва, операции сложения и
умножения, а также константы~$0$ и~$1$. Интерпретация:
комплексные числа. Предикат $x=\gamma$, где $\gamma$\т некоторое
комплексное число, выразим для рациональных $\gamma$ и невыразим
для иррациональных $\gamma$.

В самом деле, если $\gamma$ иррационально, то оно может быть
алгебраическим или трансцендентным. В первом случае рассмотрим
многочлен из $\bbQ[x]$ минимальной степени, обращающийся в~$0$ в
точке $\gamma$; по предположению он имеет степень больше~$1$ и
потому имеет другой корень~$\gamma'$. В алгебре доказывается (с
использованием трансфинитной индукции или леммы Цорна, а также
базисов трансцендентности), что существует автоморфизм $\bbC$
над $\bbQ$, переводящий $\gamma$ в $\gamma'$.

В случае трансцендентного $\gamma$ мы используем такой
факт: для любых трансцендентных $\gamma_1, \gamma_2\hm\in\bbC$
существует автоморфизм поля $\bbC$ над $\bbQ$, который
переводит~$\gamma_1$ в~$\gamma_2$.

\label{r-definability}%
Отметим, что для поля $\bbR$ вместо $\bbC$ такое рассуждение не
проходит, так как это поле не имеет нетривиальных автоморфизмов.
(Отношение порядка в нём выразимо: положительные числа суть квадраты,
поэтому любой автоморфизм сохраняет порядок. Поскольку
автоморфизм оставляет на месте все рациональные числа, он должен
быть тождественным.)

В этом случае предикат $x=\gamma$ выразим тогда и только тогда,
когда $\gamma$\т алгебраическое число. Это легко следует из
теоремы Тарского\ч Зайденберга (раздел~\ref{seidenberg-tarski},
с.~\pageref{r-definability-2}).
        %
\end{itemize}

\begin{problem}
        %
Покажите, что предикат $y=x+1$ невыразим в интерпретации $(\bbZ,
{=}, f)$, где $f$\т одноместная функция $x\hm\mapsto(x+2)$.
        %
\end{problem}

\begin{problem}
        %
Покажите, что предикат $x=2$ невыразим в множестве целых
положительных чисел с предикатами равенства и \лк $x$~делит~$y$\пк.
        %
\end{problem}

\section
[Элиминация кванторов]
{Невыразимые предикаты: элиминация кванторов}
\index{Элиминация кванторов|defin}%
        \label{quantifier-elimination}%
При всей простоте метод доказательства невыразимости с помощью
автоморфизмов страдает очевидным недостатком: очень часто
требуемого автоморфизма нет. Например, натуральные числа с
операцией прибавления единицы вообще не допускают никакого
нетривиального автоморфизма. (Тем не менее там выразимо очень
немногое, как мы вскоре увидим.) Целые числа с операцией
прибавления единицы допускают автоморфизмы (сдвиги), но эти
автоморфизмы не позволяют доказать, что отношение порядка
невыразимо (поскольку оно устойчиво относительно сдвигов).

Более прямой метод доказательства состоит в том, что мы
предъявляем некоторый класс $\mathcal{E}$
предикатов, который содержит все
выразимые предикаты и не содержит интересующего нас предиката.
При этом мы доказываем, что $\mathcal{E}$ содержит все выразимые предикаты,
так: проверяем, что $\mathcal{E}$ содержит все предикаты,
выразимые атомарными формулами, а также замк\-нут относительно
логических операций (объединение, пересечение, дополнение) и
операции проекции (соответствующей навешиванию квантора
существования; квантор всеобщности выражается через квантор
существования). Часто класс $\mathcal{E}$ совпадает с классом всех
предикатов, выразимых бескванторными формулами (иногда надо
расширить сигнатуру), и потому этот метод называют методом \лк
элиминации кванторов\пк. (Это краткое описание, возможно, станет
яснее из приводимых далее примеров.)

Начнём с такого примера.
Пусть сигнатура содержит равенство,
одноместную функцию $S$ (прибавление единицы) и константу $0$.
Носителем интерпретации будет множество $\bbZ$ целых чисел,
символы сигнатуры интерпретируются естественным образом.
В этой ситуации изоморфизмов не существует, так что предыдущий
способ доказательства невыразимости здесь неприменим.

Тем не менее класс выразимых предикатов весьма ограничен: это
предикаты, выразимые бескванторными формулами. Будем называть
две формулы (рассматриваемой нами сигнатуры) эквивалентными (в
данной интерпретации), если они выражают один и тот же предикат,
то есть истинны при одних и тех же значениях переменных.

\begin{theorem}
\index{Элиминация кванторов!в $(\bbZ,=,S,0)$|defin}%
        \label{quantifier-elimination-integers-successor}%
Для всякой формулы рассматриваемой нами сигнатуры существует
эквивалентная ей бескванторная формула.
        %
\end{theorem}

\begin{proof}
        %
Будем доказывать индукцией по построению (или, если угодно, по
длине) формулы $\varphi$ существование эквивалентной ей в
$(\bbZ,=,S,0)$
бескванторной формулы. Для удобства (чтобы рассматривать один
случай, а не два) будем считать, что наша формула может
содержать только кванторы существования, но не всеобщности. Это
законно, так как формулы $\forall \xi \,\psi$ и
$\lnot\exists\xi\,\lnot\psi$ эквивалентны.

Случай, когда $\varphi$ есть атомарная формула, очевиден\т она и
так бескванторная. Если $\varphi$ является конъюнкцией, дизъюнкцией
или импликацией двух частей, достаточно заменить каждую часть
на эквивалентную бескванторную (что можно сделать по
предположению индукции).

Единственный содержательный случай\т когда формула $\varphi$
начинается с квантора существования, то есть имеет вид $\exists
x \tau$ (пусть под квантором стоит переменная~$x$).
Мы рассуждаем по индукции, поэтому можем считать,
что формула $\tau$\т бескванторная.
Она имеет (возможно) и другие параметры, скажем,
$x_1,\dots,x_n$. Чтобы подчеркнуть это, обычно вместо~$\tau$
пишут $\tau(x,x_1,\dots,x_n)$. Нам надо найти бескванторную
формулу нашей сигнатуры, эквивалентную формуле
        $$
\exists x\, \tau(x,x_1,\dots,x_n).
        $$
Формула $\tau(x,x_1,\dots,x_n)$ представляет собой булеву
комбинацию атомарных формул. Посмотрим на те атомарные формулы,
которые содержат переменную $x$. Атомарная формула представляет
собой равенство двух термов
$S(S(\dots(S(u))\dots))=S(S(\dots(S(v))\dots))$; здесь $u$ и
$v$\т либо переменные, либо константа $0$. Если переменная $x$
входит и в левую, и в правую часть, то (в этой интерпретации)
такая атомарная формула либо всегда истинна, либо всегда ложна,
и её можно заменить на какую\д нибудь тождественно истинную или
тождественно ложную
формулу, не содержащую $x$. После этого останутся
атомарные формулы, которые можно записать как
        $$
x = t_1,\quad x = t_2,\quad\dots,\quad x = t_k.
        $$
Здесь $t_i$\т либо целая константа, либо выражение вида $x_j+c$,
где $x_j$\т какая\д то другая переменная, а $c$\т целое число.
Мы позволили себе слегка отступить от канонов, разрешив прибавлять
и вычитать целые константы вместо того, чтобы применять функцию $S$
в левой и правой частях равенства. Ясно, что это не меняет класса
выразимых формул, зато позволяет оставить $x$ в левой части, а
константу перенести в правую.

Теперь сравним формулу
        $$
\varphi=\exists x \,\tau(x,x_1,\dots,x_n)
        $$
с формулой
        $$
\tau(t_1,x_1,\dots,x_n) \lor
\tau(t_2,x_1,\dots,x_n) \lor \ldots
\lor \tau(t_k,x_1,\dots,x_n),
        $$
которую мы будет обозначать~$\varphi'$.
Формула $\varphi'$ представляет собой дизъюнкцию формул,
полученных в результате подстановки различных $t_i$ вместо $x$ в
бескванторную формулу $\tau(x,x_1,\dots,x_n)$. (После
подстановки можно вернуться к обычному виду записи формулы,
заменив прибавление констант на нужное количество применений
функции $S$ с той или другой стороны равенства.)

Очевидно, что если для каких\д то значений переменных
$x_1,\dots,x_n$ формула
$\varphi'$ истинна, то для этих значений
$x_1,\dots,x_n$ истинна и формула
$\varphi$. В самом деле, если истинен
$i$\д й член дизъюнкции, то в формуле $\varphi$ в качестве $x$
можно взять значение выражения $t_i$.

Верно ли обратное? Не обязательно. Вполне возможно, что тот $x$,
который существует и делает формулу $\varphi$ истинной,
отличается от всех $t_i$. Но мы пропустили по существу только
один случай\т все такие $x$ в некотором смысле одинаковы, так
как они делают все атомарные формулы, содержащие $x$, ложными,
поэтому всё равно, какой из таких $x$ выбрать. Отметим также,
что хотя бы один такой $x$ найдётся, поскольку $\bbZ$
бесконечно, а выражений $t_i$ лишь конечное число.

Обозначим через $\varphi''$ формулу, которая получится из~$\tau$
заменой всех атомарных формул, содержащих $x$, на тождественно
ложные формулы. Сказанное выше объясняет, почему формула $\varphi$
эквивалентна дизъюнкции $\varphi'\lor\varphi''$. Мы достигли цели
\т нашли бескванторную формулу, эквивалентную формуле $\varphi$.
        %
\end{proof}

Легко понять, что отношение порядка $x>y$ не выражается
бескванторной формулой нашей сигнатуры, поскольку такая формула
может включать лишь атомарные формулы вида $x=y+c$ и для неё
случай, когда $y$ сильно больше $x$, неотличим от случая, когда
$y$ сильно меньше $x$. Тем самым мы доказали (чего нельзя было
сделать методом автоморфизмов), что отношение $x>y$ невыразимо
(в данной интерпретации данной сигнатуры).

Немного более сложное рассуждение понадобится, если добавить
к сигнатуре отношение порядка.

\begin{theorem}
\index{Элиминация кванторов!в $(\bbZ,{=},{<}, S)$|defin}%
        \label{successor-order-elimination}%
Всякая формула в $(\bbZ,{=},{<}, S)$ (где $S$\т
функция прибавления единицы) эквивалентна некоторой
бескванторной формуле. (Как говорят, $(\bbZ,{=},{<},
S)$ \emph{допускает элиминацию кванторов}.)%
        %
\end{theorem}

\begin{proof}
        %
Полностью утверждение теоремы звучит так: для всякой формулы
сигнатуры, содержащей равенство, порядок и символ $S$, найдётся
бескванторная формула той же сигнатуры, которая эквивалентна ей
в интерпретации, где носителем является $\bbZ$, а символы
сигнатуры интерпретируются естественным образом. (В дальнейшем
мы будем опускать такие пояснения.)

Доказательство следует прежней схеме. Правда, теперь атомарных
формул больше\т помимо формул $x=t_i$ у нас будут формулы
$x<t_i$. Поэтому нельзя рассчитывать на то, что все значения
$x$, не встречающиеся среди $\{t_1,\dots,t_k\}$, ведут себя
одинаково, и наш приём с выделением случая, когда все равенства
ложны, более не проходит.

Как же быть? Для данных значений $x_1,\dots,x_n$ числа
$t_1,\dots,t_k$ делят числовую ось (точнее, множество $\bbZ$
целых чисел) на промежутки, и для выяснения истинности формулы
$\varphi$ нам надо попробовать (помимо всех $t_i$) хотя бы по
одному числу из каждого промежутка. Это будет гарантировано,
если мы напишем дизъюнкцию, в которую, помимо всех формул
$\tau(t_i,x_1,\dots,x_n)$, войдут также формулы
$\tau(t_i+1,x_1,\dots,x_n)$ и $\tau(t_i-1,x_1,\dots,x_n)$. Это
позволяет нам обойтись без формулы $\varphi''$ и благополучно
завершить доказательство.
        %
\end{proof}

\begin{problem}
        %
Проверьте, что добавление константы~$0$ к этой сигнатуре
не препятствует элиминации кванторов.
        %
\end{problem}

Что будет, если мы из этой сигнатуры удалим функцию $S$? Легко
понять, что класс выразимых множеств не изменится, так как
$y=S(x)$ можно выразить как \лк $y$ является наименьшим
элементом, б\'ольшим $x$\пк. Однако при этом мы использовали
кванторы, так что для $(\bbZ,{=},{<})$ элиминация
кванторов невозможна.

\begin{problem}
        %
Убедитесь, что  в самом деле формула $y\hm=S(x)$ не
эквивалентна никакой бескванторной формуле этой сигнатуры.
        %
\end{problem}

Часто такой переход приходится выполнять в обратном направлении:
у нас есть некоторая ситуация, в которой элиминация кванторов не
проходит. Мы обходим эту трудность, добавив некоторые
выразимые предикаты и функции в нашу сигнатуру, после чего
элиминация кванторов удаётся. В этом случае мы получаем описание
всех выразимых предикатов (предикат выразим, если он
записывается бескванторной формулой расширенной сигнатуры). Мы
встретимся с такой ситуацией дальше, говоря об арифметике
Пресбургера%
\index{Арифметика!Пресбургера}%
\glossary{\Пресбургер}
(раздел~\ref{presburger}).

\medskip

В некоторых случаях рассуждение упрощается, если
привести бескванторную формулу к дизъюнктивной
нормальной форме. Вот один из таких примеров.
        %
\begin{theorem}
\index{Элиминация кванторов!в $(\bbQ,{=},{<})$|defin}%
        \label{quantifier-elimination-dense}
Всякая формула в $(\bbQ,{=},{<})$ эквивалентна
некоторой бескванторной формуле.
        %
\end{theorem}

\begin{proof}
        %
Как всегда, достаточно рассмотреть случай формулы вида
        $$
\exists x \,\tau(x,x_1,\dots,x_n),
        $$
где $\tau(x,x_1,\dots,x_n)$\т бескванторная формула. Формулу~$\tau$
можно считать формулой в дизъюнктивной нормальной форме
(теорема~\ref{dnf-cnf}, с.~\pageref{dnf-cnf}).
Напомним, это означает, что $\tau$ представляет
собой дизъюнкцию конъюнкций, а каждая конъюнкция соединяет
несколько литералов (атомарных формул или их отрицаний).

В данном случае можно избавиться от отрицаний, заменив формулу
$\lnot(x=y)$ на $((x<y)\lor (x>y))$, а формулу $\lnot(x<y)$\т на
$((x=y)\lor(x>y))$. После этого надо воспользоваться
дистрибутивностью и вновь прийти к дизъюнктивной нормальной
форме\т с большим числом членов, но уже без отрицаний.

Теперь надо воспользоваться тем, что квантор
существования (\лк бесконечную дизъюнкцию\пк) можно переставлять с обычной дизъюнкцией.
Точнее говоря, мы пользуемся тем, что формулы
$\exists x\,(\tau_1 \lor \tau_2)$ и
$\exists x\,\tau_1 \lor  \exists x\,\tau_2$
эквивалентны. (Белый или чёрный единорог существует тогда и
только тогда, когда существует белый единорог или существует
чёрный единорог.) Это обстоятельство позволяет заменить
формулу
        $$
\exists x\, (\tau_1\lor\tau_2\lor\ldots\lor\tau_n)
        $$
на
        $$
\exists x\, \tau_1\lor\exists x\,\tau_2\lor\ldots\lor\exists x\,\tau_n
        $$
и дальше разбираться с каждой из формул поодиночке.

Итак, нам осталось преобразовать к бескванторному виду формулу
        $$
\exists x \, (\rho_1 \land \rho_2 \land\ldots\land\rho_k),
        $$
где каждая из формул $\rho_i$ соединяет какие\д то две переменные
знаком $=$ или $<$ (напомним, что от отрицаний мы уже избавились).

Некоторые из формул $\rho_i$ не содержат переменной $x$. Тогда
их можно вынести за квантор: если $x$ не является параметром
формулы $\alpha$, то формулы
%        $$
%\exists x\, (\alpha\land\beta)\quad\text{и}\quad
%\alpha \land  \exists x\,\beta
%        $$
$\exists x\, (\alpha\land\beta)$ и $\alpha \land  \exists x\,\beta$
эквивалентны
(если $\alpha$ истинно для некоторых значений параметров, то
в обеих формулах его можно опустить; если $\alpha$ ложно, то
обе формулы ложны при этих значениях параметров).

Вынеся такие формулы, можно считать, что под квантором остались
лишь формулы вида $x<x_i$, $x=x_i$ и $x>x_i$, сравнивающие
переменную $x$ с какими\д то другими переменными. Если там
есть хоть одно равенство, то квантор существования вырождается\т
его можно удалить вместе с переменной $x$, заменив её на ту
переменную, которой она равна. Например, формулу $\exists
x\,((x\hm=y)\hm\land A(x))$ можно заменить на $A(y)$.

Итак, остался случай, когда переменная $x$ встречается лишь в
неравенствах. Другими словами, нас спрашивают, найдётся ли
значение $x$, большее каких\д то переменных и меньшее
каких\д то других. Если все ограничения на $x$ одного знака
(только снизу или только сверху), то такое значение~$x$
существует при любых
значениях других переменных (поскольку в множестве $\bbQ$ нет
ни наибольшего, ни наименьшего элементов). Что делать, если
есть ограничения разных знаков? Пусть наша формула, например,
имеет вид
        $$
\exists x \, ((x>a)\land (x>b)\land(x<c)\land(x<d)).
        $$
Как записать условия на $a,b,c,d$, при которых это верно,
не используя кванторов? Надо написать такую формулу:
        $$
(a<c)\land (a<d)\land (b<c) \land (b<d).
        $$
Мы хотим написать, что наибольшая из нижних границ меньше
наименьшей из верхних, но поскольку заранее неизвестно, какая
будет наибольшей и какая наименьшей, мы пишем, что любая нижняя
граница меньше любой верхней. Поскольку множество $\bbQ$
является плотным (между любыми двумя элементами найдётся
третий), то эта формула равносильна исходной.

Так, постепенно сводя дело ко всё более простым случаям,
мы завершили рассуждение.
        %
\end{proof}

Заметим, что в этом доказательстве из свойств рациональных чисел
мы использовали лишь отсутствие наибольшего и наименьшего
элемента и плотность. Поэтому все наши преобразования остаются
эквивалентными для любого упорядоченного множества с такими
свойствами, а не только для $\bbQ$. Применив эти преобразования
к замкнутой формуле (формуле без параметров), мы получим или
тождественно истинную формулу, или тождественно ложную (только
надо добавить в язык константы для истины и лжи, чтобы не
использовать фиктивных переменных, когда надо написать
тождественно истинное или тождественно ложное выражение). Отсюда
мы заключаем, что во всех плотных упорядоченных множествах без
первого и последнего элемента справедливы одни и те же формулы
нашей сигнатуры. Как говорят, все такие
        \label{dense-sets-elementary-equivalence}%
множества \emph{элементарно эквивалентны}%
\index{Множества, элементарная эквивалентность}%
\index{Элементарная эквивалентность}
с точки
зрения нашей сигнатуры, см.~раздел~\ref{elementary-equivalence}.
(Другое доказательство этого
факта можно получить, используя теорему Левенгейма\ч Сколема%
\index{Теорема!Левенгейма\ч Сколема об элементарной подмодели}%
\glossary{\Лёвенгейм}%
\glossary{\Сколем}
о счётной подмодели и теорему об изоморфизме счётных
плотных линейно упорядоченных множеств без первого и
последнего элементов, см.~раздел~\ref{lowenheim-skolem-1}.)

В частности, мы доказали, что для рациональных и
действительных чисел истинны одни и те же формулы сигнатуры
$({=},{<})$.

Ещё одним побочным продуктом нашего рассуждения (как и других
рассуждений об элиминации кванторов) является способ выяснить,
будет ли данная замкнутая формула истинной или ложной в
рассматриваемой интерпретации. Для этого надо привести её к
бескванторному виду и посмотреть, получится ли \И\ или \Л.
Другими словами, элиминация кванторов устанавливает
\emph{разрешимость}%
\index{Разрешимая теория|defin}
элементарной теории рациональных чисел с
отношениями равенства и порядка.

\medskip

Элиминация кванторов остаётся возможной (и рассуждение даже
немного упрощается), если рациональные (или действительные)
числа рассматривать не только с равенством и порядком, но и со
сложением и рациональными константами.
        \label{q-plus-order-elimination}%
В этом случае можно
воспользоваться приведённой ранее схемой с конечным
представительным набором термов. В самом деле, пусть $x$\т
переменная, которую (вместе с квантором существования по ней) мы
хотим элиминировать. Все атомарные формулы, её содержащие, можно
\лк разрешить\пк\ относительно~$x$, получив некоторое количество
формул вида $x=t_i$, $x>t_i$ и $x<t_i$, где $t_i$\т линейные комбинации
остальных переменных с рациональными коэффициентами. (Разрешение
рациональных коэффициентов вместо целых
ничего не меняет, так как можно
привести всё к общему знаменателю и получить целые коэффициенты,
затем перенести отрицательные коэффициенты в другую часть, а
положительные заменить многократным сложением.)

Затем в качестве представительного набора надо взять набор,
состоящий, во\д первых, из всех $t_i$, во\д вторых, из всех средних
арифметических $(t_i+t_j)/2$, и, наконец, из выражений $t_i-1$ и
$t_i+1$. Ясно, что как бы ни расположились точки $t_i$ на
числовой оси, этот набор захватит как минимум по одной точке из
каждого промежутка (средние арифметические нужны для интервалов,
а прибавление и вычитание единицы\т для лучей по краям).

\begin{problem}
        %
Проведите это рассуждение подробно.
        %
\end{problem}

Возможность элиминации кванторов в только что рассмотренной ситуации
$(\bbQ, {=},
{<}, {+},$\,рациональные конс\-тан\-ты$)$ имеет
интересное геометрическое применение.

\begin{theorem}
\index{Теорема!о разбиении квадрата|defin}%
\index{Разрезание квадрата}%
\index{Квадрат, разрезание}%
        \label{square-dissected}%
Предположим, что единичный квадрат разрезан на несколько квадратов.
Тогда все они имеют рациональные стороны.
        %
\end{theorem}

\begin{proof}
        %
Пусть дано такое разрезание с $n$ меньшими квадратами.
Напишем формулу с $3n$ параметрами ($n$ из которых соответствуют
размерам меньших квадратов, а~$2n$\т координатам их левых верхних
углов), которая говорит, что эти параметры действительно задают
разрезание квадрата (квадраты содержатся внутри единичного, не
имеют общих точек и всякая точка единичного квадрата покрывается
хотя бы одним из меньших квадратов). Навесив кванторы существования
по переменным, задающим координаты, получим формулу $F(x_1,\dots,x_n)$
с параметрами $x_1,\dots,x_n$,
которая истинна, когда из квадратов размеров $x_1,\dots,x_n$ можно
составить единичный квадрат.

Элиминация кванторов позволяет считать, что формула~$F$
бескванторная, то есть представляет собой логическую комбинацию
равенств и неравенств вида $c_1x_1\hm+\ldots\hm+c_nx_n\hm+c_0\hm=0$ и
$c_1x_1\hm+\ldots\hm+c_nx_n\hm+c_0\hm> 0$ с различными рациональными
коэффициентами. Посмотрим на все выражения, стоящие в левой
части таких равенств и неравенств. Подставим в них числа
$\overline{x}_1,\dots,\overline{x}_n$, соответствующие
данному нам разрезанию. При этом получится истинная формула
$F(\overline{x}_1,\dots,\overline{x}_n)$, в которой некоторые из левых
частей равенств и неравенств
будут
равны нулю, а другие нет. Временно забудем про те,
которые не равны нулю, а равные нулю будем воспринимать как левые части
уравнений с неизвестными $x_1,\dots,x_n$ (с
нулём в правой части) независимо от того, входили ли они в $F$ как
левые части уравнений или неравенств.
Получится система уравнений, для
которой числа $\overline{x}_1,\dots,\overline{x}_n$ будут
решениями. Если эти числа являются единственными её решениями,
то они рациональны (например, потому, что правила Крамера
для решения системы уравнений содержат отношения определителей с
рациональными элементами). Покажем, что другого быть не
может.

В самом деле, если решение не единственно, то есть целая прямая,
проходящая через точку $\overline{x}=
\langle\overline{x}_1,\dots,\overline{x}_n\rangle$ и лежащая в
множестве решений системы. Все точки прямой, достаточно близкие
к $\overline{x}$, неотличимы от $\overline{x}$ с точки зрения
формулы~$F$ и потому делают формулу $F$ истинной. В самом деле,
левые части, равные нулю в $\overline{x}$, равны нулю на всей
прямой, а левые части, отличные от нуля в точке~$\overline{x}$,
сохраняют знак в некоторой окрестности этой точки. Пусть
$\langle h_1,\dots,h_n\rangle$\т направляющий вектор этой
прямой. Тогда при всех достаточно малых значениях $t$ из
квадратов размеров
        $$
\overline{x}_1+th_1,
\overline{x}_2+th_2,\dots,
\overline{x}_n+th_n
        $$
можно составить единичный
квадрат. Но это невозможно. Чтобы убедиться
в этом, достаточно оставить логику и вернуться к геометрии:
площади этих квадратов в сумме должны равняться $1$, но
площадь каждого есть многочлен второй степени по $t$,
и коэффициент при $t^2$ положителен, поэтому сумма таких
многочленов не может тождественно равняться единице ни на каком
интервале.
        %
\end{proof}

Насколько существенна в этом рассуждении ссылка на возможность
элиминации кванторов? В принципе можно было бы рассуждать так.
Пусть дано разрезание квадрата. Посмотрим на конфигурацию,
образуемую меньшими квадратами, и напишем равенства и
неравенства на размеры частей, которые гарантируют, что в этой
конфигурации нет щелей и перекрытий. (Далее продолжаем
рассуждение как раньше.) Конечно, возникает вопрос: почему мы
уверены, что такую систему уравнений и неравенств можно
написать? Глядя на конкретную конфигурацию, это сделать легко,
но как провести это рассуждение строго для общего случая, не
вполне понятно.

\medskip

Изложенные методы позволяют провести элиминацию кванторов
и описать выразимые множества во многих ситуациях; несколько
простых примеров такого рода предлагаются в качестве задач.
Два более сложных примера (арифметика Пресбургера и
теория сложения и умножения действительных чисел) разбираются
в двух следующих разделах.

\begin{problem}
        %
Опишите выразимые предикаты, проведя элиминацию
кванторов (и расширив сигнатуру, если нужно) для
       (\textbf{а})~%
$(M, {=})$, где $M$\т произвольное
бесконечное множество;
       (\textbf{б})~%
$(\bbQ, {=},{+})$;
       (\textbf{в})~%
$(\bbQ, {=}, S)$, где $S$\т операция прибавления единицы;
       (\textbf{г})~%
$(\bbN, {=}, S)$, где $S$\т операция прибавления единицы;
       (\textbf{д})~%
$(\bbN, {=}, S, P)$, где $S$\т операция прибавления единицы,
а $P$\т одноместный предикат \лк быть степенью двойки\пк.
        %
\end{problem}
